% ==============================================================================
% File: template_paper.tex
% TEMPLATE MẪU SLIDE CHO BÀI BÁO MỚI (DÙNG ĐỂ IMPORT VÀO MAIN.TEX)
% Cách dùng: 
%   1. Sao chép file này: cp template_paper.tex <ten_bai_bao>_paper.tex
%   2. Điền nội dung vào các khung bên dưới
%   3. Thêm vào main.tex: \input{<ten_bai_bao>_paper.tex}
% ==============================================================================

% Tên Section sẽ hiển thị trên mục lục và thanh tiêu đề của Beamer
\section{Tên Bài Báo Viết Tắt (PAPER-NAME)}

% ------------------------------------------------------------------------------
% Slide 0: Tổng quan bài báo
% ------------------------------------------------------------------------------
\begin{frame}{PAPER-NAME: Tổng quan bài báo}
    \begin{itemize}
        \item \textbf{Tên bài báo:} \textit{Full Title of the Research Paper}
        \item \textbf{Tác giả \& Tổ chức:} Tên tác giả chính et al. (Lab / Trường / Công ty, Năm xuất bản)
        \item \textbf{Ý tưởng cốt lõi:} Mô tả ngắn gọn trong 1-2 câu về giải pháp chính bài báo đề xuất.
        \item \textbf{Đóng góp nổi bật:}
        \begin{itemize}
            \item Đóng góp 1: Giải quyết nút thắt hoặc hạn chế của các phương pháp đi trước.
            \item Đóng góp 2: Đề xuất kiến trúc, hàm mục tiêu hoặc kỹ thuật mới.
            \item Đóng góp 3: Đạt kết quả thực nghiệm SOTA hoặc cải thiện vượt trội về tốc độ/tài nguyên.
        \end{itemize}
    \end{itemize}
\end{frame}

% ------------------------------------------------------------------------------
% Slide 1: Bối cảnh & Thách thức
% ------------------------------------------------------------------------------
\begin{frame}{1. PAPER-NAME: Bối cảnh \& Thách thức}
    \textbf{Bối cảnh nghiên cứu:}
    \begin{itemize}
        \item Nêu bài toán thực tế cần giải quyết và hướng tiếp cận truyền thống (ví dụ: autoregressive, standard diffusion).
        \item Tại sao hướng tiếp cận này quan trọng trong thời điểm hiện tại.
    \end{itemize}
    \vspace{0.2cm}
    \textbf{Hạn chế của các phương pháp hiện tại (Pain Points):}
    \begin{itemize}
        \item Điểm yếu 1: Chi phí tính toán cao, tốc độ suy luận chậm hoặc yêu cầu nhiều tài nguyên.
        \item Điểm yếu 2: Lỗi xấp xỉ lý thuyết hoặc suy giảm chất lượng trong các tình huống thực tế.
        \item Hậu quả cụ thể: Mô hình bị sụp đổ, sinh kết quả kém hoặc khó mở rộng quy mô.
    \end{itemize}
\end{frame}

% ------------------------------------------------------------------------------
% Slide 2: Ý tưởng cốt lõi & Khung làm việc (Framework)
% ------------------------------------------------------------------------------
\begin{frame}{2. PAPER-NAME: Khung làm việc (Framework)}
    \begin{itemize}
        \item \textbf{Đột phá tư duy (Key Insight):} Điểm khác biệt căn bản trong cách tiếp cận bài toán.
        \item \textbf{Quy trình tổng thể (Overall Pipeline):}
        \begin{itemize}
            \item \textit{Đầu vào / Biểu diễn:} Không gian dữ liệu đầu vào (ví dụ: latent embedding, unconstrained space).
            \item \textit{Tiến trình biến đổi:} Cách thức xử lý trung gian (ví dụ: continuous flow, recurrent state).
            \item \textit{Đầu ra / Giải mã:} Cách chuyển đổi về kết quả cuối cùng.
        \end{itemize}
        \item \textbf{Ưu thế so với thiết kế cũ:} Loại bỏ các bước xấp xỉ dư thừa hoặc tối ưu hóa đường đi của mô hình.
    \end{itemize}
\end{frame}

% ------------------------------------------------------------------------------
% Slide 3: Phương pháp luận & Công thức toán then chốt
% ------------------------------------------------------------------------------
\begin{frame}{3. PAPER-NAME: Phương pháp luận toán học}
    \begin{itemize}
        \item \textbf{Định nghĩa biến \& Quá trình:} Mô tả công thức nền tảng:
        $$z_t = \alpha_t x_1 + \beta_t x_0, \quad t \in [0, 1]$$
        \item \textbf{Phương trình vi phân / Toán tử ánh xạ:}
        $$\dot{x}_t = v_t(x_t), \quad \text{hoặc} \quad X_{s,t}(x_s) = x_t$$
        \item \textbf{Tính chất lý thuyết cốt lõi:}
        \begin{itemize}
            \item Tính chất 1: Bảo toàn không gian nghiệm hoặc nằm trọn trên miền xác suất hợp lệ.
            \item Tính chất 2: Thỏa mãn điều kiện giải tích (ví dụ: semigroup property, Lipschitz continuity).
        \end{itemize}
    \end{itemize}
\end{frame}

% ------------------------------------------------------------------------------
% Slide 4: Hàm mục tiêu & Cơ chế huấn luyện
% ------------------------------------------------------------------------------
\begin{frame}{4. PAPER-NAME: Hàm mục tiêu \& Huấn luyện}
    \begin{itemize}
        \item \textbf{Hàm mất mát chính (Loss Function):}
        $$\mathcal{L}_{\text{total}} = \mathcal{L}_{\text{main}} + \lambda \mathcal{L}_{\text{aux}}$$
        \item \textbf{Phân tích thành phần hàm loss:}
        \begin{itemize}
            \item $\mathcal{L}_{\text{main}}$: Đảm bảo mô hình khớp với dữ liệu hoặc neo vào mô hình giáo viên.
            \item $\mathcal{L}_{\text{aux}}$: Ràng buộc điều hòa hoặc chưng cất bước nhảy.
        \end{itemize}
        \item \textbf{Chiến lược lấy mẫu & Ổn định huấn luyện:}
        \begin{itemize}
            \item Phân bổ thời gian (time schedule / reparameterization).
            \item Kỹ thuật lấy mẫu biên hoặc warm-up tham số.
        \end{itemize}
    \end{itemize}
\end{frame}

% ------------------------------------------------------------------------------
% Slide 5: Cơ chế suy luận & Dẫn hướng (Inference & Guidance)
% ------------------------------------------------------------------------------
\begin{frame}{5. PAPER-NAME: Cơ chế suy luận (Inference)}
    \begin{itemize}
        \item \textbf{Thuật toán lấy mẫu (Sampling Algorithm):}
        \begin{itemize}
            \item Chế độ sinh ít bước (few-step / one-step) hoặc tích phân số nhiều bước (ODE/SDE).
            \item Công thức cập nhật trạng thái tại mỗi bước $t_n \to t_{n+1}$.
        \end{itemize}
        \item \textbf{Cơ chế dẫn hướng (Inference Guidance):}
        \begin{itemize}
            \item \textit{Tự dẫn hướng (Autoguidance / CFG):} Tăng chất lượng mẫu không điều kiện.
            \item \textit{Dẫn hướng theo phần thưởng (Reward Guidance):} Lái mẫu sinh theo thuộc tính mong muốn (sentiment, topic, logic constraints).
        \end{itemize}
    \end{itemize}
\end{frame}

% ------------------------------------------------------------------------------
% Slide 6: Kết quả thực nghiệm
% ------------------------------------------------------------------------------
\begin{frame}{6. PAPER-NAME: Kết quả thực nghiệm}
    \textbf{Đánh giá định lượng so với các Baselines:}
    \begin{itemize}
        \item \textbf{Chất lượng mẫu (Perplexity / Accuracy / FID):}
        \begin{itemize}
            \item Benchmark 1: PAPER-NAME đạt \textbf{X.XX} vs Baseline 1 (Y.YY), Baseline 2 (Z.ZZ).
            \item Benchmark 2: Duy trì độ đa dạng cao, không bị sụp đổ mẫu (mode collapse).
        \end{itemize}
        \item \textbf{Tốc độ \& Hiệu năng tính toán:}
        \begin{itemize}
            \item Tăng tốc độ suy luận gấp \textbf{k lần} so với các phương pháp trước.
            \item Tiết kiệm số bước lấy mẫu và chi phí FLOPs.
        \end{itemize}
    \end{itemize}
\end{frame}

% ------------------------------------------------------------------------------
% Slide 7: Nghiên cứu triệt tiêu & Kết luận
% ------------------------------------------------------------------------------
\begin{frame}{7. PAPER-NAME: Phân tích triệt tiêu \& Kết luận}
    \textbf{Ablation Study (Kiểm chứng các quyết định thiết kế):}
    \begin{itemize}
        \item Yếu tố 1 (Hàm loss / Tham số hóa): Khẳng định lựa chọn A vượt trội hơn lựa chọn B.
        \item Yếu tố 2 (Time schedule / Lấy mẫu): Chứng minh tầm quan trọng của kỹ thuật đề xuất.
    \end{itemize}
    \vspace{0.2cm}
    \textbf{Kết luận \& Hạn chế:}
    \begin{itemize}
        \item \textbf{Ưu điểm chính:} Đóng góp một giải pháp mới, hiệu quả và có cơ sở lý thuyết vững chắc.
        \item \textbf{Hạn chế & Hướng tương lai:} Chi phí bộ nhớ hoặc hướng mở rộng cho các bài toán quy mô lớn hơn.
    \end{itemize}
\end{frame}

